% ===================================================================
%  بخش ۳ — روش‌شناسی و مدل شبیه‌سازی
%  منبع: متن موجود مقاله، بدون تغییر محتوایی (فقط شماره‌ی بخش عوض شد:
%  ۲ → ۳)، مطابق جدول نگاشت در evidence/integration_plan.md
%
%  هندسه‌ی زیر، هندسه‌ی همان اجرایی است که اعداد بخش ۴ از آن آمده‌اند و
%  از پارامترهای آبجکت‌های ساخته‌شده در Lumerical استخراج شده است:
%    استوانه: radius = 10nm ، z span = 40nm ، مرکز z = 0
%    دو کره‌ی کلاهک: radius = 10nm در z = +20nm و z = -20nm
%    ⇒ L_total = 40 + 10 + 10 = 60nm ، D = 20nm ، نسبت منظری = 3
%    راس کلاهک بالا: z = 30nm ؛ دوقطبی: z = 35nm (گپ 5nm)
%    Mesh Override: dx=dy=dz = 2nm روی ناحیه‌ی 40×40×80 nm
%    trans_box = 400nm ، ناحیه‌ی FDTD = 600nm با PML
%  ثبت راستی‌آزمایی: evidence/source_audit.md، بخش «هندسه و قله‌های
%  راستی‌آزمایی‌شده». اگر اجرای نویی با هندسه‌ی دیگر انجام شد، هم این
%  اعداد و هم simulation/nanorod_purcell.lsf را هم‌زمان اصلاح کنید.
% ===================================================================
\section{روش‌شناسی و مدل شبیه‌سازی}
\label{sec:method}

شبیه‌سازی‌های این پژوهش با استفاده از روش تفاضل محدود در حوزه زمان
سه‌بعدی (3D-FDTD) انجام شده است \cite{taflove2005computational}. ساختار
مورد مطالعه شامل یک نانومیله طلا با طول کل $L = \nm{60}$ و قطر
$D = \nm{20}$ (نسبت منظری ۳) است که از یک بدنه‌ی استوانه‌ای به طول
$\nm{40}$ و دو کلاهک نیم‌کروی به شعاع $R = D/2 = \nm{10}$ در دو انتها
تشکیل شده است؛ بنابراین راس کلاهک بالایی در $z = \nm{30}$ روی محور تقارن
قرار می‌گیرد. داده‌های تجربی
ضریب شکست مختلط طلا از مدل جانسون و کریستی
\cite{johnson1972optical} اقتباس شده و با مدل تحلیلی چندجمله‌ای
درون‌نرم‌افزاری بر ساختار فیت گردیده است. محیط پیرامون به عنوان
دی‌الکتریک همگن (هوا با ضریب شکست $n=1$) در نظر گرفته شد.

گسیل‌گر نوری به صورت یک چشمه دوقطبی الکتریکی نقطه‌ای (Electric Dipole) با
راستای قطبش موازی با محور طولی نانومیله ($z$) و در فاصله
$d = \nm{5}$ از راس کلاهک، یعنی در $z = \nm{35}$، قرار داده شده است.
ناحیه محاسباتی مکعبی با ضلع $\nm{600}$ توسط
لایه‌های جاذب کاملاً منطبق (PML) احاطه شده تا از هرگونه بازتاب غیرفیزیکی
ممانعت به عمل آید. جهت دستیابی به دقت بالا و حل صحیح گرادیان شدید میدان در
گپ و مرز مشترک فلز-دی‌الکتریک، یک شبکه مش محاسباتی موضعی (Mesh Override)
با گام $dx = dy = dz = \nm{2}$ بر ناحیه‌ای به ابعاد
$40 \times 40 \times \nm{80}$ که نانومیله و محل چشمه را می‌پوشاند تعریف
گردید.

اعتبار توصیف الکترودینامیکی موضعی کلاسیک در این فاصله بررسی شد. فاصله‌ی
۵ نانومتری گسیل‌گر از سطح فلز حدود یک مرتبه بزرگی بالاتر از گاف‌های
زیرنانومتری است که در آن‌ها تونل‌زنی الکترونی آغاز می‌شود؛ در گاف‌های
بزرگ‌تر از حدود ۱ نانومتر، تقویت میدان با کاهش گاف همچنان رشد می‌کند و
تنها در زیر این مقیاس است که تونل‌زنی سبب افت شدید آن می‌گردد. قطر
$\nm{20}$ نانومیله نیز بالاتر از حد حدوداً ۱۰ نانومتری است که تقریب پاسخ
موضعی برای نانوساختارهای فلزی معتبر شمرده می‌شود \cite{xu2023quantum}.
با این حال تصریح می‌شود که پاسخ غیرموضعی (پاشندگی فضایی) در مقیاس‌های
چندنانومتری می‌تواند تصحیح‌های کمّی بر جای بگذارد
\cite{babicheva2023optical}؛ اعداد این پژوهش در چارچوب پاسخ موضعی
کلاسیک و بر پایه‌ی داده‌های تجربی جانسون و کریستی گزارش می‌شوند.

شار توان تابش‌شده به میدان دوردست از طریق انتگرال‌گیری عددی بردار
پوئینتینگ بر روی سطوح یک جعبه بسته شش‌وجهی فرضی
($\text{trans\_box}$) به ضلع $\nm{400}$ پیرامون ساختار و دوقطبی محاسبه
شده است. این جعبه از یک سو کل نانومیله و چشمه را در بر می‌گیرد و از سوی
دیگر حداقل $\nm{100}$ از مرزهای PML فاصله دارد تا بازتاب عددی مرزها بر
شار توان اثر نگذارد:
%
\begin{equation}
  T(\lambda) = \frac{P_{\text{rad}}}{P_0}
  = \frac{\iint_{\text{box}} \mathbf{S} \cdot d\mathbf{A}}{P_0}.
  \label{eq:T}
\end{equation}
%
همچنین فاکتور پورسل کل ($\Fp$) که بیانگر نسبت نرخ واپاشی کل در حضور
نانوآنتن ($\Gamma_{\text{tot}}$) به نرخ واپاشی در فضای همگن
($\Gamma_0$) است، مستقیماً از کار کل انجام‌شده توسط میدان واکنشی بر روی
خود چشمه دوقطبی به دست آمد:
%
\begin{equation}
  \Fp(\lambda) = \frac{\Gamma_{\text{tot}}}{\Gamma_0}
  = \frac{P_{\text{tot}}}{P_0}.
  \label{eq:Fp}
\end{equation}
%
بازده کوانتومی ظاهری و میزان توان اتلافی غیرتابشی به فرم حرارتی (تلفات
اهمی) نیز بر پایه روابط زیر محاسبه می‌شوند:
%
\begin{align}
  \eta_a(\lambda) &= \frac{P_{\text{rad}}}{P_{\text{tot}}}
  = \frac{T(\lambda)}{\Fp(\lambda)}, \label{eq:eta} \\[2pt]
  \frac{P_{\text{loss}}}{P_0} &= \Fp(\lambda) - T(\lambda).
  \label{eq:loss}
\end{align}
%
تفکیک صریح رابطه‌های~\eqref{eq:T} و~\eqref{eq:Fp} — یعنی محاسبه‌ی
مستقل توان تابشی و توان کل، به‌جای استخراج یکی از دیگری با یک ضریب
ثابت — پیش‌شرط تحلیلی است که در زیربخش~\ref{sec:results-implication}
ارائه می‌شود.
